\scanpage{119}
\mspage{125}
$c_{\deleted{\unsure{1|2}}}'$) Toute $\unsure{\mathbb D|D}$-équivalence est \deleted{\unsure{une}} $1$-complète

$c''$) Toute $\mathbb D$-équivalence est $2$-complète.

\uline{NB} On a

\diagram{\begin{tikzcd}[ampersand replacement=\&,row sep=small,column sep=small]
\& c''\arrow[d,Rightarrow] \&\\
\& c'\arrow[dl,Rightarrow]\arrow[dr,Rightarrow] \&\\
c_1 \& \& c_2
\end{tikzcd}}

\unsure{tautologiquement}. Je ne connais que HOT pour satisfaire \deleted{\unsure{les}} \inserted{\unsure{à l'une|l'une} des deux} plus faibles conditions $c_1,c_2$. Et en fait, sous réserve d'écrire la théorie du dérivateur HOT, il doit être vrai qu'elle satisfait même la forme la plus forte $c''$. (Et les HOT-équiv doivent être simplement les Hot-équivalences.) Cf le ``lemme'' p.~\unsure{104|109}.)

La proposition suivante est une tautologie

\uline{Proposition 10} Supposons que $\mathbb D$ satisfasse la condition Der 3 $c_1$ (p.~124), p. ex. $\mathbb D=\mathrm{HOT}$. Alors \unsure{la|les} notion de $\mathbb D$-fibration coïncide avec celle de $\mathbb D$-fibration cohomologique. Elle est donc stable par composition. \inserted{Pour que $f$\inserted{$:X\unsure{\mathord{-}|\to}Y$} soit une $\mathbb D$-fibration il \unsure{f. et s.} que pour $Y''\xrightarrow{p}Y'$ sur $Y$, si $p$ est une $\mathbb D$-équiv., il en est de même pour $q\unsure{=|:}X''\to X'$.}

\uline{Corollaire 1} Si $f:X\to Y$ dans Diag est $\mathbb D$-propre \deleted{(resp. $\mathbb D$-lisse)}, alors $f$ est une $\mathbb D$-fibration cohomologique \unsure{sss|ssi} \deleted{pour} toute flèche $y'-y$ dans $Y$, $X_{/y'}-X_{/y}$ est une $\mathbb D$-équiv. (i.e. sss $f$ est une $\mathbb D$-fibration \inserted{coh.} faible).

\unsure{Appliquons} le critère c) de prop. \unsure{6|5}.
% archive sheet number: 119
